% 第九卷 6.5.2 节：候选树与树注意力掩码。
% 节点 1 周五、2 周六（深度 1）；3 下午、4 上午为 1 的子节点，5 下午、6 上午为 2 的子节点（深度 2）。
% 掩码：节点 i 可见节点 j 当且仅当 j = i 或 j 是 i 的祖先；前缀对所有节点可见，图中省略。
\begin{tikzpicture}[
  x=1cm,y=1cm,
  every node/.style={font=\figurefont\scriptsize,text=FigureInk},
  nd/.style={rectangle,rounded corners=1.8pt,draw=FigureBlue,line width=0.6pt,fill=FigureBlueFill,minimum height=5mm,minimum width=11mm,inner sep=1pt},
  rt/.style={nd,draw=FigureStroke,fill=FigurePaper},
  idx/.style={font=\figurefont\scriptsize,text=FigureMuted},
  seen/.style={draw=white,line width=0.5pt,fill=FigureBlue!70},
  hidden/.style={draw=FigureGrid,line width=0.4pt,fill=white},
  hd/.style={font=\figurefont\bfseries\footnotesize,text=FigureInk}
]
  % 候选树。
  \node[hd] at (1.7,1.85) {候选树};
  \node[rt] (r) at (0,0) {当前前缀};
  \node[nd] (n1) at (1.75,0.8) {1 周五};
  \node[nd] (n2) at (1.75,-0.8) {2 周六};
  \node[nd] (n3) at (3.6,1.2) {3 下午};
  \node[nd] (n4) at (3.6,0.4) {4 上午};
  \node[nd] (n5) at (3.6,-0.4) {5 下午};
  \node[nd] (n6) at (3.6,-1.2) {6 上午};
  \foreach \a/\b in {r/n1,r/n2,n1/n3,n1/n4,n2/n5,n2/n6} {
    \draw[-{Stealth[length=1.6mm]},FigureStroke,line width=0.6pt] (\a.east) -- (\b.west);
  }
  \node[idx] at (1.75,-1.75) {深度1};
  \node[idx] at (3.6,-1.75) {深度2};

  % 掩码矩阵。
  \def\s{0.42}
  \def\xo{5.6}
  \def\yo{1.3}
  \node[hd] at ({\xo+3*\s},1.85) {树注意力掩码};
  \foreach \i in {1,...,6} {
    \foreach \j in {1,...,6} {
      \pgfmathtruncatemacro{\pa}{\i<=2 ? 0 : (\i<=4 ? 1 : 2)}
      \pgfmathtruncatemacro{\vis}{(\j==\i || \j==\pa) ? 1 : 0}
      \ifnum\vis=1
        \path[seen] ({\xo+(\j-1)*\s},{\yo-(\i-1)*\s}) rectangle ++(\s,-\s);
      \else
        \path[hidden] ({\xo+(\j-1)*\s},{\yo-(\i-1)*\s}) rectangle ++(\s,-\s);
      \fi
    }
    \node[idx,anchor=east] at ({\xo-0.05},{\yo-(\i-0.5)*\s}) {\i};
    \node[idx,anchor=south] at ({\xo+(\i-0.5)*\s},{\yo+0.02}) {\i};
  }
  \draw[FigureStroke,line width=0.6pt] (\xo,\yo) rectangle ++({6*\s},{-6*\s});
  \node[idx,anchor=north] at ({\xo+3*\s},{\yo-6*\s-0.08}) {列：键节点；行：查询节点};
  \path[seen] ({\xo+6*\s+0.4},0.5) rectangle ++(0.25,-0.25);
  \node[idx,anchor=west] at ({\xo+6*\s+0.7},0.375) {可见};
  \path[hidden] ({\xo+6*\s+0.4},0.05) rectangle ++(0.25,-0.25);
  \node[idx,anchor=west] at ({\xo+6*\s+0.7},-0.075) {不可见};
\end{tikzpicture}
